\section{Introduction}
\label{sec:intro}
A brain--computer interface (BCI) turns measured brain activity into commands
for a device~\cite{wolpaw2002}. In a motor-imagery (MI) BCI the user imagines
a movement, which weakens the mu and beta rhythms over the motor
cortex~\cite{pfurtscheller1999,pfurtscheller2001}, and a decoder learns to
recognise this change in the electroencephalogram (EEG). MI-BCIs are studied
for stroke rehabilitation~\cite{cervera2018,biasiucci2018} and for motor
augmentation, in which people learn to control an extra robotic
digit~\cite{kieliba2021,eden2022}. Wang et al.\ joined the two: healthy
people~\cite{wang2024jbhi} and then stroke patients~\cite{wang2024tnsre}
imagined flexing a robotic sixth finger worn on the wrist. Their public
dataset of \nSubjects{} patients now opens this paradigm to independent
study~\cite{wang2026dataset}.

Every MI-BCI must first be calibrated: the user performs labelled trials so
that the decoder can learn their patterns~\cite{lotte2015}. For a patient,
each minute of calibration is effort spent before any therapy starts.
Transfer learning reuses data from other people and shortens calibration for
healthy users~\cite{zanini2018,he2020,wu2022review}. Whether it helps stroke
patients learn a command with no natural counterpart, such as a sixth
finger, is untested.

Such a study needs an evaluation free of leakage first. The dataset paper
reports 85--86\,\% for telling imagined hand from imagined sixth-finger
movements, with spatial filters fitted on all trials of a session before
cross-validation~\cite{wang2026dataset}, a common source of optimistic
estimates~\cite{kapoor2023,brookshire2024}. We re-evaluate the dataset with
every step fitted on training data only, and base every later number on
that evaluation.

We ask how much calibration a stroke patient needs before a decoder
recognises rest, the affected hand and the sixth finger, and whether data
from other patients reduce it. Our contributions are:
\begin{enumerate}
  \item a leakage-free benchmark of this sixth-finger stroke dataset, with an
        integrity audit that found duplicated recordings;
  \item to our knowledge, the first calibration study of a supernumerary
        (novel) motor command in stroke patients, with pre-registered
        hypotheses, for a new user and for a new command whose hand trials
        are already recorded;
  \item a comparison of Riemannian and deep transfer for this population,
        with stroke-specific analyses (hemisphere mirroring, other days) as
        exploratory work; and
  \item a demonstrator that replays held-out recordings through a decoder and
        moves a virtual hand with a sixth finger.
\end{enumerate}
